\section{SYSTEM DESIGN AND IMPLEMENTATION}

\subsection{Architecture Overview}
CurricuPath is a web-based academic advisor for Wits EFA03 and EFA04
students. Language-model components help interpret wording, while
deterministic curriculum rules make eligibility decisions. A request passes
through message interpretation, catalogue resolution, rule evaluation,
response checking and presentation.

The system returns an eligibility decision when evidence is sufficient, asks
for clarification when a required student fact is missing, or refers the case
for manual academic review when evidence is incomplete, conflicting or
requires discretion. The back end is a Python FastAPI service and the user
interface is a React application.

\subsection{Data and Evidence}
CSV files provide curriculum modules, rules, source records, aliases and
module outlines, derived from the 2026 EBE handbook \cite{wits2026ebe} and
School registration workbooks. Each rule stores its source and page
reference, and each source record carries a verification status that shows
whether it can support an eligibility decision. Curriculum requirements are
stored as nested all-of and any-of expressions, so each condition can be
checked and cited separately.

At run time the catalogue is held in PostgreSQL; if the database is
unavailable, the service falls back to the validated CSV files. Session
information, such as profile facts and pending clarification questions, is
temporary and expires with the session.

\subsection{Request Interpretation}
A deterministic parser identifies the request type, module, programme and
year from module codes, names, programme abbreviations and controlled
aliases. It also reads statuses the student states, such as ``I failed
ELEN3009A'' or ``I am taking Control I'', which a language model cannot set.
In the deployed advisor, Gemini (\texttt{gemini-3.5-flash-lite}) reads each
message but may only supply a request type the parser could not identify; it
cannot introduce a module code or programme the student did not name. In the
reported evaluation, Gemini was used only as a fallback for unresolved
messages.

Ambiguous module names or a missing programme lead to a clarification
question. While a question is open, a short reply is treated as its answer
rather than as a new request. For unresolved messages, spelling correction
and catalogue matching are applied before Gemini suggests a module or
programme, and every suggestion is checked against the catalogue. If Gemini
is unavailable, the deterministic path continues and a warning is recorded.

\subsection{Eligibility and Clarification}
Eligibility is evaluated through ordered checks for source verification,
programme membership, prerequisites, corequisites, exclusions and timetable
clashes. Each condition is marked satisfied, unsatisfied or unknown, and an
unknown condition is never treated as satisfied. The rule engine returns
\textit{eligible}, \textit{ineligible}, \textit{needs clarification} or
\textit{manual review required}, together with the relevant rule and source.

Two checks are enabled by configuration and were enabled for the
evaluation: School registration-grid requirements, which take priority over
handbook requirements for the same module and rule type, and timetable
clashes, which use pairs listed in the School workbooks and apply to
repeated, current or planned modules. When a required fact is missing, one
question covers every unknown module in the failed rule; answers persist
for the session, and further questions are asked only while each reply adds
information. Declared concessions, transfer credits and similar
circumstances are referred for manual academic review.

\subsection{Response Checks and Technology Selection}
Before a response is shown, deterministic checks confirm that the decision,
module codes and cited rules match the rule-engine result; a failing response
is designed to be replaced by a fixed template. Course recommendations show
each option's eligibility status and unmet requirements, so ineligible
options are labelled rather than hidden.

Python reads the School workbooks directly and keeps the rule engine, data
import and evaluation scripts in one environment. FastAPI's typed request
models reject malformed profiles before they reach the rule engine and
generate API documentation \cite{fastapiDocs}. React holds the profile,
pending questions and evidence views on the client \cite{reactDocs}.
PostgreSQL stores linked curriculum, rule and source records, mapped by
SQLAlchemy and versioned with Alembic \cite{sqlalchemyDocs,alembicDocs}.
pgvector keeps Hugging Face embeddings next to the records they describe for
similarity search \cite{pgvectorDocs,huggingFaceEmbeddings}; these support
retrieval and recommendation, never eligibility. Gemini returns structured
output \cite{googleStructuredOutput}, and a request may lower, but never
raise, the server's language-model setting.
